% Chapter 13: Subject-Verb Agreement

\chapter{Subject-Verb Agreement: Making Verbs Match}

\section{Pedagogical Purpose}
Subject-verb agreement is a fundamental rule for grammatical accuracy in English. After learning about subjects, nouns, pronouns, and verbs individually, students are now ready to understand how these elements must work together harmoniously. This chapter is crucial for developing correct sentence construction and editing skills.

\begin{learningobjectives}
The learner can:
\begin{itemize}
    \item identify the subject and verb in a simple sentence.
    \item explain the rule of subject-verb agreement.
    \item use singular verbs with singular subjects.
    \item use plural verbs with plural subjects.
    \item recognise the special agreement pattern for "I" and "you."
    \item correct common subject-verb agreement errors, including sentences with words between the subject and the verb.
\end{itemize}
\end{learningobjectives}

\section{Prerequisite Check}
\begin{quickcheck}
Underline the subject and circle the verb in each sentence.
\begin{enumerate}
    \item \underline{Rohan} \textcircled{plays} football.
    \item \underline{Maya and Leo} \textcircled{are} friends.
\end{enumerate}
\end{quickcheck}

\begin{rememberbox}
You already know that every sentence needs a subject and a verb (Chapter 1), and that verbs show action or a state of being (Chapter 8). Now you will learn that the subject and the verb must also \textbf{agree} with each other in number.
\end{rememberbox}

\section{Chapter Opener}
Rohan and Maya are talking about their favourite animals, and Ms. Sharma notices something interesting about the way they speak.

\textbf{Rohan:} "My favourite animal \textbf{is} a lion."

\textbf{Maya:} "Oh, I like lions too! But my favourite animals \textbf{are} elephants."

\textbf{Ms. Sharma:} "Interesting! Rohan, you said 'lion \textbf{is},' and Maya, you said 'elephants \textbf{are}.' Why did you use two different words for 'to be'?"

\textbf{Rohan:} "Because 'lion' is just one, and 'elephants' are many?"

\textbf{Ms. Sharma:} "Exactly, Rohan! You've noticed something very important about how English works. Let's explore it properly."

\section{Core Explanation}
\textbf{Subject-Verb Agreement} means that the verb in a sentence must match its subject in number.
\begin{itemize}
    \item If the subject is \textbf{singular} (one), the verb must be \textbf{singular}.
    \item If the subject is \textbf{plural} (more than one), the verb must be \textbf{plural}.
\end{itemize}

\begin{examplebox}
    The \textbf{cat} \textbf{sleeps} all day. (Singular subject "cat" with singular verb "sleeps")
    The \textbf{cats} \textbf{sleep} all day. (Plural subject "cats" with plural verb "sleep")
\end{examplebox}

\subsection*{A. The Basic Rule}
Think of it like this: Singular subjects like to have verbs that end in \textbf{-s}. Plural subjects like verbs without the \textbf{-s}.

\begin{table}[h!]
\centering
\begin{tabular}{|l|l|l|}
\hline
\textbf{Subject} & \textbf{Verb Form} & \textbf{Example} \\
\hline
\textbf{Singular} (one) & Verb + \textbf{s} & The \textbf{bird} sing\textbf{s}. \\
\textbf{Plural} (many) & Verb (base form) & The \textbf{birds} sing. \\
\hline
\end{tabular}
\end{table}

\subsection*{B. Agreement with 'be', 'have', and 'do'}
Some of the most common verbs have special forms for singular and plural subjects.

\begin{table}[h!]
\centering
\begin{tabular}{|l|l|l|}
\hline
\textbf{Verb} & \textbf{Singular Subject} & \textbf{Plural Subject} \\
\hline
\textbf{be} & is, was & are, were \\
\textbf{have} & has & have \\
\textbf{do} & does & do \\
\hline
\end{tabular}
\end{table}

\subsection*{C. Special Subjects: 'I' and 'You'}
The pronouns \textbf{I} and \textbf{you} are special.
\begin{itemize}
    \item \textbf{I} (which is singular) uses a plural verb form.
    \begin{itemize}
        \item I \textbf{sing}. (not I sings)
        \item I \textbf{have} a pen. (not I has)
    \end{itemize}
    \item \textbf{You} can be singular or plural, but it always uses a plural verb form.
    \begin{itemize}
        \item You \textbf{are} my friend.
        \item You \textbf{sing} well.
    \end{itemize}
\end{itemize}

\begin{warningbox}
    \textbf{Words Between the Subject and the Verb}
    Sometimes, a phrase comes between the subject and the verb. Don't be tricked! The verb must agree with the real subject, not the words in between.
    \begin{itemize}
        \item The \textbf{basket} of apples \textbf{is} on the table. (The subject is \textit{basket}, not \textit{apples}.)
        \item The \textbf{boy} with the red shoes \textbf{runs} very fast. (The subject is \textit{boy}, not \textit{shoes}.)
    \end{itemize}
\end{warningbox}

\section{Summary}
\begin{summarybox}
    \begin{itemize}
        \item The \textbf{subject} and \textbf{verb} in a sentence must agree in number.
        \item A \textbf{singular subject} (one) needs a \textbf{singular verb}.
        \item A \textbf{plural subject} (more than one) needs a \textbf{plural verb}.
        \item "I" and "you" are special — they always use the plural-style verb form.
        \item A phrase between the subject and the verb (like "of apples") does not change what the true subject is.
        \item Collective nouns (like "team" or "class") are usually treated as singular.
    \end{itemize}
\end{summarybox}